\documentclass{beamer}
\usepackage[utf8]{inputenc}
\usepackage[T1]{fontenc}
\usepackage[spanish]{babel}
\usepackage{amsmath,amssymb,bm}

% Apariencia similar a las diapositivas originales
\usetheme{default}
\usecolortheme{rose}
\setbeamertemplate{navigation symbols}{}
\setbeamertemplate{blocks}[rounded][shadow=true]
\setbeamercolor{block title}{fg=blue!70!black,bg=blue!15}
\setbeamercolor{block body}{bg=blue!7}
\setbeamercolor{frametitle}{fg=blue!70!black}
\setbeamercolor{title}{fg=blue!70!black}
\setbeamerfont{title}{series=\bfseries}
\setbeamerfont{frametitle}{series=\bfseries}
\setbeamertemplate{footline}{%
  \hfill\insertframenumber/10\hspace{1em}\vspace{0.5em}}

\title{SME0823 Modelos Lineales Generalizados:\\ Estimación Puntual Frecuentista}
\author{\textbf{Osvaldo Anacleto Jr}}
\institute{\textbf{ICMC-USP}}
\date{16 de agosto de 2018}

\begin{document}

% ---------- Página 1 ----------
\begin{frame}[plain]
  \titlepage
\end{frame}

% ---------- Página 2 ----------
\begin{frame}
  \begin{block}{Estructura de la clase}
    \begin{itemize}
      \item Repaso de inferencia\\
            (archivo \texttt{GLMaula4bRevisaoInferencia.pdf})
      \vspace{0.8em}
      \item Función de verosimilitud de un MLG
      \vspace{0.8em}
      \item Definición del problema y obtención de las ecuaciones de score
      \vspace{0.8em}
      \item Matriz de información de Fisher para MLGs
    \end{itemize}
  \end{block}

  \vspace{2em}
  {\color{blue}\bfseries Referencias}\\
  Cordeiro \& Demétrio (2008, Cap.~3,2), Paula (2004, caps.~1.5 y 1.6)
  y Wakefield (2013, cap.~6.5)
\end{frame}

% ---------- Página 3 ----------
\begin{frame}
  \colorbox{violet!60!black}{%
    \parbox{\dimexpr\linewidth-2\fboxsep}{\color{white}
      {\color{yellow!80!orange}\bfseries Objetivo:} dado un MLG, obtener
      estimaciones puntuales del vector de parámetros
      $\bm{\beta}=(\beta_0,\dots,\beta_p)^{\mathsf{T}}$ del componente
      sistemático del modelo}}

  \vspace{1em}
  Estas estimaciones se obtendrán mediante el método de máxima
  verosimilitud.

  \vspace{0.8em}
  Sea $Y$ una variable aleatoria referente a la variable respuesta (etiqueta)
  del MLG, con distribución perteneciente a la familia exponencial.

  \vspace{0.8em}
  La función de verosimilitud del modelo se escribe con base en la función
  de densidad de $Y$,
  \begin{equation}
    f(y;\theta,\phi)=\exp\Big\{\phi^{-1}\big[y\theta-b(\theta)\big]+c(y,\phi)\Big\}
  \end{equation}

  \vspace{0.5em}
  con $\mathrm{E}[Y]=\mu=b'(\theta)$ y $\mathrm{Var}[Y]=\phi\,\mathrm{V}(\mu)$,
  donde $\mathrm{V}(\mu)=b''(\theta)$ es la función de varianza.
\end{frame}

% ---------- Página 4 ----------
\begin{frame}{Función de verosimilitud}
  Sean $y_1,y_2,\dots,y_n$ una muestra de $n$ observaciones independientes
  de la variable respuesta $Y$. La función de verosimilitud está dada por
  {\footnotesize
  \[
    \mathrm{L}(\bm{\theta};y_1,\dots,y_n,\phi)
    =\prod_{i=1}^{n} f(y_i;\theta_i,\phi)
    =\prod_{i=1}^{n}\exp\Big\{\phi^{-1}\big[y_i\theta_i-b(\theta_i)\big]+c(y_i,\phi)\Big\}
  \]}
  donde $\bm{\theta}^{\mathsf{T}}=(\theta_1,\dots,\theta_n)$ son los
  parámetros canónicos.

  \vspace{1em}
  Desde el punto de vista computacional y de análisis, es conveniente
  considerar la función de log-verosimilitud
  {\footnotesize
  \[
    \mathrm{l}(\bm{\theta};\phi,\bm{y})
    =\sum_{i=1}^{n}\mathrm{l}_i
    =\sum_{i=1}^{n}\log\big(\mathrm{L}(\theta_i;\phi,y_i)\big)
    =\sum_{i=1}^{n}\Big[\phi^{-1}\big[y_i\theta_i-b(\theta_i)\big]+c(y_i,\phi)\Big]
  \]}
\end{frame}

\end{document}
